\subsubsection{文件读写}
在src文件夹中的file.c文件中实现文件的读写、删除等函数,其实方法类似任务一,只是在查找目录项、FAT表时为了确保互斥访问和一致性,需要前后申请信号量保护。
重点修改在删除文件时需要确保没有读者或写者在使用该文件,否则无法删除,具体代码如下:
\begin{ccode}{删除文件函数}
bool fs_delete(const char *filename)
{
    pthread_mutex_lock(&fs_mutex);

    for (int i = 0; i < MAX_ENTRY; i++) {
        if (Directory[i].ac == 1 &&
            Directory[i].type == 0 &&
            strcmp(Directory[i].name, filename) == 0) {

            // 检查文件是否正被打开：
            //  reader_count > 0 说明有读者
            //  sem_trywait(&rw_lock) 失败说明有写者
            pthread_mutex_lock(&Directory[i].mutex);
            if (Directory[i].reader_count > 0) {
                pthread_mutex_unlock(&Directory[i].mutex);
                pthread_mutex_unlock(&fs_mutex);
                printf("Failed to delete file: %s is in use\n", filename);
                return false;
            }
            pthread_mutex_unlock(&Directory[i].mutex);

            if (sem_trywait(&Directory[i].rw_lock) != 0) {
                pthread_mutex_unlock(&fs_mutex);
                printf("Failed to delete file: %s is in use\n", filename);
                return false;
            }
            sem_post(&Directory[i].rw_lock);  // 成功拿到写锁，立即释放 

            int cur = Directory[i].start_block;
            while (cur != -1) {
                int next = FAT[cur];
                FAT[cur] = 0;
                if (next == -1) break;
                cur = next;
            }

            Directory[i].ac = 0;
            pthread_mutex_unlock(&fs_mutex);
            printf("File deleted successfully: %s\n", filename);
            return true;
        }
    }

    pthread_mutex_unlock(&fs_mutex);
    printf("Failed to delete file: %s\n", filename);
    return false;
}
\end{ccode}
如果检查到reader\_count大于0,说明还有读者在访问文件,无法删除;如果检查到sem\_trywait(\&rw\_lock)返回失败,说明有写者在访问文件,同样无法删除。只有当文件空闲没有进程在访问时,才能删除文件,具体删除操作于任务一相同,释放占有的FAT表项和Directory目录项即可。